%&pdflatex
%%%%%%%%%%%%%%%%%%%%%%%%%%%%%%%%%%%%%%%%%%%%%%%%%%%%%%%%%%%%%%%%%%%%%%%%%%%
% DICCIONARIO EN CASTELLANO PARA \TeX,
% por
% JUAN LUIS VARONA,
% escrito para el
% Primer Encuentro del Grupo de Usuarios de TeX Hispanohablantes,
% EGUTH'99, Madrid, 13 y 14 de septiembre 1999.
%     FECHA COMIENZO ESCRITURA: 12 de julio de 1999.
%     FECHA \'ULTIMA MODIFICACI\'ON: 1 de diciembre de 1999.
%     A\~NADIDA REFERENCIA A SU PUBLICACI\'ON: 26 DE AGOSTO DE 2004.
% Escrito en LaTeX2e, por Juan Luis Varona (para problemas del texto,
% contactar con jvarona@dmc.unirioja.es).
%%%%%%%%%%%%%%%%%%%%%%%%%%%%%%%%%%%%%%%%%%%%%%%%%%%%%%%%%%%%%%%%%%%%%%%%%%%


\documentclass[a4paper,12pt]{article}

% Para que la footnote con la referncia no a\~nada una p\'agina:
\addtolength{\textheight}{0.5cm}

\usepackage[spanish]{babel}

\usepackage{hyperref}
% Si no se usa hyperref, usar esto:
%\usepackage{url}
% SE PUEDE COGER EN:
% ftp://CTAN/tex-archive/macros/latex/contrib/other/misc/url.sty
% SE USA PARA LOGRAR CORTES DE PALABRA EN PAGINAS WEB Y SIMILARES.

% SI NO SE CARGA hyperref, INCLUIR ESTO:
%\def\href#1#2{\url{#2}}

\date{}

\title{DICCIONARIO EN CASTELLANO\\ PARA \TeX}


\author{Juan Luis Varona%
\renewcommand\thefootnote{}
\renewcommand\footnotemark{}
\footnote{Esta ponencia fue presentada en el
Primer Encuentro del Grupo de Usuarios de TeX Hispanohablantes,
EGUTH'99, celebrado en la Universidad Polt\'ecnica de Madrid
los d\'{\i}as 13 y 14 de septiembre 1999.
Fue editado en las Actas del Encuentro (distribuidas a los
asistentes mediante fotocopias), y posteriormente publicado en
\textit{TeXemplares} \textbf{3} (2001), n.\textsuperscript{o}~2,
32--36.}
}

\begin{document}

\maketitle

\begin{abstract}
Se presenta un diccionario electr\'onico en castellano destinado a
ser usado, fundamentalmente, por verificadores ortogr\'aficos de \TeX.

El diccionario es una actualizaci\'on del que estaba disponible
hasta ahora y que ya se emplea, por ejemplo, en Excalibur para
Macintosh, en WinEdt para Windows, y en amSpell para MS-DOS.
Consiste en una recopilaci\'on de casi 300 mil palabras.
\end{abstract}

\section{Verificadores ortogr\'aficos}

Cuando escribimos un texto, sea del tipo que sea, es bastante com\'un que, por
mucho cuidado que pongamos, cometamos alguna errata. No s\'olo es posible que
nos equivoquemos al recordar la ortograf\'{\i}a de una palabra, sino que
nuestra impericia escribiendo nos lleva a menudo a cambiar unas letras por otras
al teclear. Para intentar minimizar esto en lo posible, unos programas llamados
verificadores ortogr\'aficos vienen en nuestra ayuda.

Un verificador ortogr\'afico comprueba que cada palabra de un texto es una de
las palabras posibles en el idioma. Obviamente, un verificador nunca
descubrir\'a un error cuando, al escribir mal una palabra, hemos escrito otra
tambi\'en existente, incluso aunque no tenga ning\'un sentido gramatical en la
frase en que aparece.

En lo que a nosotros ---usuarios de \TeX--- respecta, es importante tener en
cuenta que un verificador ortogr\'afico que quiera ser \'util para comprobar
documentos escritos en \TeX{} debe comprender, al menos m\'{\i}nimamente, la
sintaxis del lenguaje: de lo contrario, marcar\'a como errores los comandos de
\TeX, las f\'ormulas, etc. As\'{\i} pues, deberemos utilizar un verificador
ortogr\'afico especializado en \TeX.

Por otra parte, el diccionario de palabras admisibles puede estar construido
de dos maneras, lo que da lugar a dos tipos de verificadores ortogr\'aficos. La
primera consiste simplemente en una enumeraci\'on de todas las palabras
posibles. As\'{\i}, cada palabra debe aparecer con todas sus variaciones
permitidas al a\~nadirle distintos prefijos y sufijos. Por ejemplo, los
sustantivos y adjetivos deben aparecer en masculino, femenino, singular y
plural; y los verbos conjugados.

La segunda no es una simple lista de palabras con todas sus derivaciones
posibles. Aqu\'{\i}, pod\'{\i}amos pensar que en el diccionario tenemos s\'olo
la ra\'{\i}z de la palabra y, adem\'as, unas reglas de derivaci\'on. Por ejemplo,
para cada sustantivo deben aparecer las reglas de c\'omo formar masculinos,
femeninos y plurales; y cada verbo debe ir asociado con sus reglas de
conjugaci\'on.

Programas del primer tipo son Excalibur para Macintosh, amSpell para PC con
MS-DOS y WinEdt para Windows. Del segundo tipo, podemos destacar ispell, que
est\'a disponible para equipos unix.

En general, los del primer tipo son m\'as sencillos, pero requieren mayores
ficheros de palabras. Los del segundo tipo requieren un trabajo de dise\~no y
programaci\'on m\'as importante, pero pueden llegar a ser m\'as
efectivos. Por ejemplo, para incluir aumentativos y diminutivos s\'olo hace
falta a\~nadir unas pocas reglas. En uno de los del primer tipo requiere
incluir en la lista cada palabra en aumentativo y diminutivo.

Por otra parte, unos idiomas se adaptan mejor que otros a utilizar verificadores
de uno u otro tipo. Por ejemplo, en ingl\'es no hay femeninos ni plurales de
adjetivos, y los verbos se conjugan muy poco. As\'{\i}, no merece mucho la pena
el esfuerzo de programaci\'on que se necesita para utilizar verificadores del
segundo tipo. El espa\~nol, en cambio, tiene gran cantidad de formas verbales
que, en un verificador del primer tipo, es necesario incluir individualmente.
Otros idiomas, como el alem\'an y el vasco, aun tienen m\'as mecanismos de
formaci\'on de palabras, como la declinaci\'on y la yuxtaposici\'on
(que puede llegar a ser bastante com\'un), lo que hace que el n\'umero de
palabras permitidas crezca desmesuradamente. Para estos idiomas, una estrategia
del segundo tipo parece ser bastante m\'as recomendable. El castellano
podr\'{\i}a ser un caso intermedio.

\section{La lista de palabras en castellano}

En lo que a nosotros respecta, nos vamos a ocupar aqu\'{\i} \'unicamente de
verificadores ortogr\'aficos del primer tipo, de los basados en una lista de
palabras. Desde hace tiempo, existen diversos programas que pueden ser usados
para comprobar la correcci\'on ortogr\'afica de textos escritos en \TeX, unos
gratis y otros no excesivamente caros. Pero no hab\'{\i}a ninguna lista de
palabras en espa\~nol suficientemente completa y de distribuci\'on gratuita.
Obviamente, s\'{\i} que exist\'{\i}an implementaciones en programas comerciales
pero que, en la pr\'actica, no serv\'{\i}an para nada si lo que quer\'{\i}amos
era comprobar un documento \TeX. As\'{\i} pues, hace ya a\~nos, el que escribe
se plante\'o elaborar una lista lo bastante amplia como para que pudiera ser
usada con comodidad.

Ahora se presenta una actualizaci\'on de la lista, que contiene m\'as de 40.000
palabras nuevas. Concretamente, la lista actual consta de 292.342 palabras,
mientras que la anterior ten\'{\i}a 247.049. Es importante hacer notar que
incluye multitud de palabras relacionadas con la matem\'atica, f\'{\i}sica e
inform\'atica, que son las disciplinas donde m\'as se utiliza \TeX. Muchas de
ellas no suelen estar presentes en diccionarios de car\'acter general.

Es de destacar que el usuario tiene acceso al diccionario fuente y puede
modificarlo a su antojo para su uso personal. Aunque ser\'{\i}a deseable que la
versi\'on que se distribuye en general permanezca unificada. As\'{\i}, las
futuras mejoras y a\~nadidos de palabras se hacen siempre sobre la misma lista.
El mecanismo para indicar la versi\'on de la lista es simple: basta indicar el
n\'umero de palabras que contiene.

\section{Implementaciones para diversos sistemas operativos}

La lista de palabras que aqu\'{\i} se presenta puede usarse con los siguientes
verificadores ortogr\'aficos:

\begin{subsection}{Excalibur}

Es un programa para Macintosh, desarrollado por Robert Gottshall y Rick Zaccone.
Es gratuito. Su p\'agina web es
\url{http://www.eg.bucknell.edu/~excalibr/excalibur.html}, y puede conseguirse
tambi\'en, mediante ftp an\'onimo, en \url{ftp://ftp.eg.bucknell.edu/pub/mac/}
(en el subdirectorio \url{Excalibur-dictionaries/} se pueden encontrar
diccionarios en multitud de idiomas). En CTAN, se encuentra en el directorio
\url{tex-archive/systems/mac/support/excalibur/}.

Una caracter\'{\i}stica de Excalibur resulta muy \'util a los escritores en
\TeX. Una palabra que contenga letras no inglesas (acentos o e\~nes, en el caso
del espa\~nol) puede escribirse con notaci\'on de 7 bits o de 8 bits. No es este
lugar para comentar las ventajas y desventajas que tiene uno y otro m\'etodo,
puesto que cualquier usuario medianamente avanzado las conoce. Con Excalibur,
basta tener la lista de palabras escrita con notaci\'on de 8 bits.
As\'{\i}, si una palabra de un texto est\'a escrita con notaci\'on de 7 bits,
Excalibur se encarga de traducir los 7 bits a 8 bits (en cualquier
combinaci\'on entendible por \TeX) antes de comparar con las palabras de la
lista.

\end{subsection}

\begin{subsection}{WinEdt}

Es un potente y vers\'atil editor destinado a la creaci\'on de documentos en
\TeX{} que funciona en Windows. Entre otras cosas, tiene un verificador
ortogr\'afico. Su autor es Aleksander Simonic y es un programa
\textit{shareware}. El precio de una licencia personal para uso educativo es
de 40 d\'olares. Su p\'agina web es \url{http://home.istar.ca/~winedt} o
\url{http://www.winedt.com}. En CTAN, puede encontrarse en el directorio
\url{tex-archive/systems/win32/winedt/}. Asimismo, en el subdirectorio
\url{dict/} se pueden encontrar diccionarios en bastantes idiomas.

\end{subsection}

\begin{subsection}{amSpell}

Es un verificador ortogr\'afico para documentos \TeX{}
que funciona bajo MS-DOS. Su autor es A.~Merckens y es gratuito. Puede
descargarse, en CTAN, en el directorio \url{tex-archive/support/amspell/},
aunque all\'{\i} s\'olo aparece el diccionario ingl\'es. El diccionario en
espa\~nol, que es una adaptaci\'on del que aqu\'{\i} se presenta realizada por
Agust\'{\i}n Mart\'{\i}n Domingo, puede conseguirse por ftp an\'onimo en
\url{ftp://tex.unirioja.es/pub/tex/dict-pc/}.

\end{subsection}
\bigskip

Para finalizar, \'unicamente a\~nadir que todos estos programas, y las
\'ultimas versiones del diccionario espa\~nol adaptado a ellos, pueden
obtenerse en la Universidad de La Rioja (de donde procede el que escribe esta
nota), mediante ftp an\'onimo, en
\url{ftp://tex.unirioja.es/pub/tex/}. Una vez ah\'{\i}, los directorios
\url{dict-mac/}, \url{dict-win/} y \url{dict-pc/} contienen,
respectivamente, los diccionarios para Macintosh (Excalibur), Windows (WinEdt)
y MS-DOS (amSpell) aqu\'{\i} descritos.

\bigskip

\noindent\textbf{Agradecimientos:}
Gracias a todos los que han colaborado notific\'andome erratas de la versi\'on
anterior de la lista de palabras, o envi\'andome nuevas palabras para incluir.
En particular, a H\'ector M.~Monacci y a Josu Zabaleta.

\bigskip

\begin{flushright}
{\small Juan Luis Varona,\\
Departamento de Matem\'aticas y Computaci\'on,
Universidad de La Rioja,\\
C/ Luis de Ulloa s/n, 26004 Logro\~no, ESPA\~NA.\\
e-mail: \href{mailto:jvarona@dmc.unirioja.es}{jvarona@dmc.unirioja.es}\\
URL: \url{http://www.unirioja.es/dptos/dmc/jvarona/welcome.html}}
\end{flushright}

\end{document}

